\chapter{Conclusão}
\label{cap:conclusao}

O projeto cumpre o objetivo de integrar Prolog num sistema pericial interativo e
explicável. A contribuição técnica principal é o motor de regras de produção
implementado em \cod{engine.pl}. A agenda determinística, a refração, a
proveniência e o índice de dependências tornam explícitos mecanismos que, noutros
motores, permanecem internos.

A base de conhecimento apresenta dimensão e variedade suficientes para
demonstrar raciocínio reativo, temporal, negativo, difuso e orientado por
objetivos. A integração com Java preserva todas as propostas e separa os canais
de controlo, enquanto a interface permite observar factos, regras e decisões. Os
testes de equivalência com Drools e a execução integral de 598 testes sem falhas
fornecem evidência de correção no contexto do simulador.

A solução deve ser entendida como um protótipo académico de apoio à decisão.
Dentro desse âmbito, oferece uma resposta completa aos critérios de avaliação:
base de conhecimento estruturada, novo motor de inferência, explicações
\emph{How}, \emph{Why} e \emph{Why not}, metaconhecimento, arquitetura
cliente-servidor e demonstração prática.
